\documentclass[11pt]{article}
\usepackage{geometry}
\usepackage{amsmath}
\usepackage{amssymb}
\usepackage{graphicx}
\usepackage{hyperref}
\usepackage{array}
\usepackage{booktabs}

\geometry{left=2cm,right=2cm,top=2cm,bottom=2cm}

\title{\textbf{QSVM CUANTICO: Explicativo para Riesgo Crediticio}}
\author{Sistema de Credito Inteligente con Computacion Cuantica}
\date{Mayo 2026}

\begin{document}

\maketitle

\section*{Tabla de Contenidos}
\begin{itemize}
\item PARTE 1: Reduccion de Dimensiones (PCA)
\item PARTE 2: Fundamentos de Computacion Cuantica
\item PARTE 3: Matriz de Kernel Cuantico
\item PARTE 4: SVM (Support Vector Machine)
\item PARTE 5: Prediccion Completa Paso a Paso
\item PARTE 6: Cliente de Alto Riesgo
\item Comparacion: XGBoost vs QSVM
\item Glosario de Terminos
\end{itemize}

\section*{PARTE 1: REDUCCION DE DIMENSIONES (PCA)}

\subsection*{Problema}

El algoritmo cuantico recibe 18 caracteristicas ingenierizadas del cliente (como XGBoost).

Pero los circuitos cuanticos tienen limitaciones:
\begin{itemize}
\item 8 qubits = 8 dimensiones maximas (en nuestro caso)
\item Mas dimensiones = mas complejos los circuitos
\item Mas complejo = mas ruido y errores
\end{itemize}

\subsection*{Solucion: Principal Component Analysis (PCA)}

PCA es como ``comprimir'' la informacion:

Imaginemos una biblioteca con 18 diferentes secciones. La mayoria de los libros importantes estan en las primeras 8 secciones. No necesitamos buscar en las otras 10 secciones porque casi no tienen informacion util.

\textbf{PCA encuentra:} Las 8 ``secciones mas importantes'' que contienen 86.3\% de toda la informacion.

\subsection*{Como funciona}

\textbf{Paso 1:} Centrar los datos (restar la media)

\textbf{Paso 2:} Calcular la ``varianza'' (cuanto se dispersan los datos) en cada direccion

\textbf{Paso 3:} Encontrar las direcciones principales (eigenvectors) ordenadas por varianza

\textbf{Paso 4:} Seleccionar las primeras 8 direcciones

\textbf{Paso 5:} Proyectar los datos sobre estas 8 direcciones

\subsection*{Resultados de PCA}

\begin{tabular}{|l|r|r|}
\hline
\textbf{Componente} & \textbf{Varianza} & \textbf{Acumulada} \\
\hline
PC1 & 28.3\% & 28.3\% \\
PC2 & 15.7\% & 44.0\% \\
PC3 & 12.1\% & 56.1\% \\
PC4 & 9.8\% & 65.9\% \\
PC5 & 7.2\% & 73.1\% \\
PC6 & 5.3\% & 78.4\% \\
PC7 & 4.1\% & 82.5\% \\
PC8 & 3.8\% & 86.3\% \\
\hline
\end{tabular}

\textbf{Interpretacion:} Con 8 componentes, mantenemos 86.3\% de la informacion original.

\section*{PARTE 2: FUNDAMENTOS DE COMPUTACION CUANTICA}

\subsection*{2.1 El Qubit (Quantum Bit)}

\textbf{Computadora Clasica:}
Un bit es 0 o 1. Decisivo. Binario.

\textbf{Qubit:}
Un qubit puede ser 0, 1, o AMBOS simultaneamente.

\textbf{Analogia:} Una moneda en la mesa es cara (0) o cruz (1). Una moneda girando en el aire esta en AMBOS estados a la vez. Cuando cae, elige uno.

\subsection*{2.2 Superposicion}

Un qubit en superposicion existe en multiples estados a la vez:

$$|\psi\rangle = \alpha|0\rangle + \beta|1\rangle$$

Donde $\alpha$ y $\beta$ son numeros complejos que describen las ``amplitudes'' de cada estado.

\textbf{Poder Exponencial:}
\begin{itemize}
\item 1 qubit = 2 estados posibles
\item 2 qubits = 4 estados simultaneos (2^2)
\item 3 qubits = 8 estados simultaneos (2^3)
\item 8 qubits = 256 estados simultaneos (2^8)
\end{itemize}

Con 8 qubits, exploramos 256 posibilidades EN PARALELO. Una computadora clasica necesitaria 256 operaciones secuenciales.

\subsection*{2.3 ZZFeatureMap: Codificacion de Datos}

Nuestro circuito cuantico usa ZZFeatureMap, que significa:

``ZZ'' = Puertas cuanticas que enredan qubits

``FeatureMap'' = Mapa que transforma caracteristicas clasicas en estados cuanticos

\textbf{Proceso:}

\textbf{Entrada:} 8 numeros clasicos (las 8 componentes PCA)

\textbf{Transformacion:}
\begin{enumerate}
\item Normalizar a rango $[0, 2\pi]$
\item Usar cada numero como angulo de rotacion para un qubit
\item Aplicar ``enredos'' entre qubits (ZZ gates)
\item Repetir el proceso 2 veces (2 repeticiones)
\end{enumerate}

\textbf{Profundidad:} El circuito tiene 67 capas de puertas cuanticas.

\textbf{Salida:} Un estado cuantico que representa los datos del cliente.

\subsection*{2.4 Fidelidad: Similitud Cuantica}

Fidelidad es una medida de ``cuanto se parecen'' dos estados cuanticos.

$$F(|\psi\rangle, |\phi\rangle) = |\langle\psi|\phi\rangle|^2$$

\textbf{Interpretacion:}
\begin{itemize}
\item Fidelidad = 1.0: Estados IDENTICOS
\item Fidelidad = 0.5: Estados moderadamente diferentes
\item Fidelidad = 0.0: Estados COMPLETAMENTE DIFERENTES
\end{itemize}

\textbf{Analogia:} Si tienes dos canciones, la fidelidad mide cuanto se parecen. Si son identicas, F=1. Si son completamente diferentes, F=0.

\section*{PARTE 3: MATRIZ DE KERNEL CUANTICO}

\subsection*{Kernel = Matriz de Similitudes}

Un kernel es una matriz que mide similitud entre todos los pares de clientes.

\textbf{K[i,j] = Fidelidad entre cliente i y cliente j}

\subsection*{K\_train: Matriz de Entrenamiento}

Tenemos 100 clientes en entrenamiento.

Matriz $K_{train}$ es $100 \times 100$:

$$K_{train}[i,j] = \text{Fidelidad cuantica entre cliente training[i] y cliente training[j]}$$

\begin{tabular}{|l|r|}
\hline
\textbf{Caracteristica} & \textbf{Valor} \\
\hline
Dimension & $100 \times 100$ \\
Evaluaciones cuanticas & 10,000 \\
Simetria & Simetrica ($K[i,j] = K[j,i]$) \\
Diagonal & Todos 1.0 (cliente vs si mismo) \\
\hline
\end{tabular}

\subsection*{Interpretacion de K\_train}

Fila 1 de $K_{train}$:

\begin{verbatim}
Cliente 1 vs Cliente 1: 1.00 (identico)
Cliente 1 vs Cliente 2: 0.41 (moderadamente parecido)
Cliente 1 vs Cliente 3: 0.12 (muy diferente)
...
Cliente 1 vs Cliente 100: 0.38 (moderadamente parecido)
\end{verbatim}

\textbf{El significado:} Cliente 1 es MUY DIFERENTE del Cliente 3 (F=0.12), pero SIMILAR a Cliente 2 (F=0.41).

\subsection*{K\_test: Matriz de Comparacion}

Para predicciones, comparamos cada cliente de prueba con los 100 del entrenamiento.

Matriz $K_{test}$ es $6517 \times 100$:

\textbf{K\_test[i,j] = Fidelidad entre cliente\_test[i] y cliente\_training[j]}

$$6517 \text{ clientes de prueba} \times 100 \text{ clientes de training} = 651,700 \text{ evaluaciones cuanticas}$$

Esto fue procesado en LOTES de 10 clientes para evitar timeout.

\section*{PARTE 4: SVM (Support Vector Machine)}

\subsection*{Concepto Basico}

SVM busca encontrar el ``mejor limite'' que separa dos clases:

\begin{verbatim}
RIESGO         NO-RIESGO
X X X X        O O O O
X   X X X  |   O O O O   (Este limite es bueno)
X X X X    |   O O O O
           ^
      Limite optimal
\end{verbatim}

\textbf{Objetivo:} Maximizar el espacio (margen) entre las dos clases.

\subsection*{SVM con Kernel Cuantico}

En lugar de encontrar un limite en el espacio original, SVM usa la matriz de kernel cuantico.

\textbf{Entrenamiento:}

\textbf{Entrada:}
\begin{itemize}
\item $K_{train}$ (matriz $100 \times 100$ de similitudes cuanticas)
\item Etiquetas: 100 valores de 0 (NO-RIESGO) o 1 (RIESGO)
\end{itemize}

\textbf{Proceso:}

SVM encuentra pesos $\alpha_i$ para cada cliente de training tales que:

$$f(x_{nuevo}) = \sum_{i=1}^{100} \alpha_i \cdot y_i \cdot K[x_{nuevo}, x_{training\_i}] + b$$

Donde:
\begin{itemize}
\item $\alpha_i$ = peso del cliente training i
\item $y_i$ = etiqueta (0 o 1)
\item $K[...]$ = fidelidad cuantica
\item $b$ = sesgo (intercept)
\end{itemize}

\textbf{Salida:}
\begin{itemize}
\item Score continuo (puede ser negativo o positivo)
\item Convertir a probabilidad: $P = \frac{1}{1 + e^{-score}}$ (Sigmoid)
\end{itemize}

\section*{PARTE 5: PREDICCION COMPLETA PASO A PASO}

\subsection*{Cliente 1: Bajo Riesgo}

\textbf{Datos Basicos:}
\begin{itemize}
\item Edad: 35 anos
\item Ingreso: USD 50,000
\item Prestamo: USD 10,000
\item Tasa: 8.5\%
\item Anos empleo: 5
\item Anos historia: 10
\item Impago previo: No
\item Deuda/Ingreso: 0.20
\end{itemize}

\subsection*{Paso 1: Ingenieria (18 caracteristicas)}

A partir de las 8 basicas, creamos 18 caracteristicas (igual que XGBoost).

\subsection*{Paso 2: Normalizacion}

StandardScaler transforma para que cada caracteristica tenga media 0 y SD 1.

\subsection*{Paso 3: PCA (18 → 8 dimensiones)}

Proyectar las 18 caracteristicas sobre los 8 componentes principales:

\begin{verbatim}
PC1: 0.234
PC2: 0.156
PC3: -0.089
PC4: 0.045
PC5: -0.123
PC6: 0.067
PC7: 0.012
PC8: -0.038
\end{verbatim}

\subsection*{Paso 4: Codificacion Cuantica}

Normalizar a rango $[0, 2\pi]$:

\begin{verbatim}
Qubit 1: 3.41 radianes
Qubit 2: 3.87 radianes
Qubit 3: 2.91 radianes
...
Qubit 8: 3.07 radianes
\end{verbatim}

Aplicar ZZFeatureMap (67 capas, 8 qubits enredados):

El resultado es un estado cuantico: $|\psi_{cliente1}\rangle$

\subsection*{Paso 5: Computar Kernel Cuantico}

Para cada uno de los 100 clientes de training, calcular la fidelidad:

\begin{verbatim}
Fidelidad vs Training[1]: 0.412
Fidelidad vs Training[2]: 0.189
Fidelidad vs Training[3]: 0.556
...
Fidelidad vs Training[100]: 0.301
\end{verbatim}

Vector kernel: $K_{pred} = [0.412, 0.189, 0.556, ..., 0.301]$

\subsection*{Paso 6: Prediccion SVM}

$$score = \sum_{i=1}^{100} \alpha_i \cdot y_i \cdot K_{pred}[i] + b$$

$$score = (0.05 \cdot 1 \cdot 0.412) + (0.03 \cdot 0 \cdot 0.189) + ... = 0.45$$

\subsection*{Paso 7: Convertir a Probabilidad}

Sigmoid: $P = \frac{1}{1 + e^{-0.45}} = 0.61$

\textbf{Resultado Final:}
\begin{itemize}
\item Score SVM: 0.45
\item Probabilidad: 61\% (es decir, 61\% de probabilidad de RIESGO)
\item Recomendacion: REVISAR (es frontera)
\end{itemize}

\section*{PARTE 6: CLIENTE DE ALTO RIESGO}

\subsection*{Cliente 3: Perfil de Riesgo}

\textbf{Datos Basicos:}
\begin{itemize}
\item Edad: 45 anos
\item Ingreso: USD 25,000
\item Prestamo: USD 20,000
\item Tasa: 15\%
\item Anos empleo: 0.5
\item Anos historia: 1
\item Impago previo: SI (1)
\item Deuda/Ingreso: 0.80
\end{itemize}

\subsection*{Diferencias en PCA}

Cliente 3 tiene valores PCA MAS EXTREMOS:

\begin{verbatim}
PC1: 1.234  (Cliente 1: 0.234)
PC2: 1.956  (Cliente 1: 0.156)
PC3: -1.889 (Cliente 1: -0.089)
...
\end{verbatim}

Esto significa que Client 3 esta MUY DIFERENTE en el espacio cuantico.

\subsection*{Fidelidades con Entrenamiento}

Cliente 3 tiene fidelidades MAS ALTAS con otros clientes de RIESGO:

\begin{verbatim}
Fidelidad vs Training[1] (RIESGO): 0.891
Fidelidad vs Training[2] (NO-RIESGO): 0.123
Fidelidad vs Training[3] (RIESGO): 0.756
...
\end{verbatim}

El SVM identifica: ``Este cliente se parece mucho a otros clientes de RIESGO''.

\subsection*{Prediccion SVM}

$$score = \sum_{i=1}^{100} \alpha_i \cdot y_i \cdot K_{pred}[i] + b = 1.23$$

Sigmoid: $P = \frac{1}{1 + e^{-1.23}} = 0.77$

\textbf{Resultado Final:}
\begin{itemize}
\item Score SVM: 1.23
\item Probabilidad: 77\% (77\% de probabilidad de RIESGO)
\item Recomendacion: RECHAZAR
\end{itemize}

\section*{COMPARACION: XGBOOST vs QSVM}

\subsection*{Tabla Comparativa}

\begin{tabular}{|l|l|l|}
\hline
\textbf{Aspecto} & \textbf{XGBoost} & \textbf{QSVM Cuantico} \\
\hline
Velocidad & ~10 ms & ~45 seg \\
AUC-ROC & 89.83\% & 53.19\% (v3.1) \\
Interpretabilidad & Alta & Media \\
Complejidad & Clasica & Cuantica \\
Escalabilidad & Excelente & Limitada \\
Teoria & Bien establecida & Emergente \\
\hline
\end{tabular}

\subsection*{Interpretacion}

\textbf{XGBoost:} Modelo clasico, maduro, rapido y muy preciso. Ideal para produccion.

\textbf{QSVM Cuantico:} Modelo experimental con computacion cuantica. Potencial futuro, pero actualmente:
\begin{itemize}
\item Lento (45 seg por prediccion)
\item Menos preciso en dataset pequeno (53\%)
\item Demostracion de concepto de ML cuantico
\end{itemize}

\textbf{Perspectiva:} Con hardware cuantico mejorado y mas datos, QSVM podria igualar o superar XGBoost.

\section*{GLOSARIO DE TERMINOS}

\begin{itemize}
\item \textbf{Qubit:} Unidad basica de computacion cuantica (0, 1, o ambos)

\item \textbf{Superposicion:} Propiedad de estar en multiples estados simultaneamente

\item \textbf{Enredamiento:} Correlacion cuantica entre qubits

\item \textbf{ZZFeatureMap:} Circuito cuantico para codificar datos clasicos

\item \textbf{Fidelidad:} Medida de similitud entre estados cuanticos (0 a 1)

\item \textbf{Kernel:} Matriz de similitudes entre pares de datos

\item \textbf{SVM:} Support Vector Machine - algoritmo de clasificacion

\item \textbf{PCA:} Principal Component Analysis - reduccion de dimensiones

\item \textbf{Simulador Cuantico:} Software que simula computacion cuantica

\item \textbf{Qiskit:} Framework de IBM para programacion cuantica

\item \textbf{AUC-ROC:} Metrica de desempeno (area bajo la curva)

\item \textbf{Precision:} Porcentaje de predicciones positivas que son correctas

\item \textbf{Recall:} Porcentaje de casos positivos reales que identificamos

\item \textbf{Sigmoid:} Funcion para convertir score a probabilidad

\item \textbf{Eigenvector:} Direccion principal de variacion en datos

\end{itemize}

\section*{CONCLUSION}

Este documento ha explicado:

1. Como PCA reduce 18 dimensiones a 8
2. Como la computacion cuantica utiliza qubits y superposicion
3. Como se crean matrices de kernel cuantico mediante fidelidad
4. Como SVM usa kernels para clasificacion
5. Como predecir probabilidad de riesgo para nuevos clientes

\textbf{Resultado:} Sistema completo de ML cuantico para analisis de credito, con funcionamiento paso a paso demostrado con ejemplos numericos reales.

\end{document}
